\documentclass[11pt,a4paper]{article}
\usepackage{../shared/luxcover}
\usepackage[utf8]{inputenc}
\usepackage[T1]{fontenc}
\usepackage[margin=1in]{geometry}
\usepackage{amsmath,amssymb,amsthm}
\usepackage{booktabs}
\usepackage{tabularx}
\usepackage{xcolor}
\usepackage{hyperref}
\hypersetup{colorlinks=true,linkcolor=black,citecolor=blue,urlcolor=blue}
\usepackage{enumitem}
\usepackage{siunitx}
\sisetup{group-separator={,},group-minimum-digits=4}
\newtheorem{proposition}{Proposition}
\newtheorem{definition}{Definition}
\newcommand{\code}[1]{\texttt{#1}}
\tolerance=800
\emergencystretch=3em

\title{Choosing the Vehicle:\\Taking an AI Equity Strategy On-Chain}

\author{
Zach Kelling\thanks{Corresponding author: research@lux.network} \\
\textit{Lux Industries Inc} \\
\textit{A subsidiary of Hanzo Industries Inc} \\
\texttt{https://lux.network}
}

\date{2026}

\begin{document}
\luxcoverpage

\begin{abstract}
There are two live precedents for putting a fund on a blockchain and they are
architecturally opposite. One is a registered closed-end interval fund whose
token enforces transfer restrictions at the contract level. The other is a family
of tokenized equities that are plain transferable tokens with compliance applied
only at issuance and redemption. The first is offerable to United States retail
investors and cannot touch decentralized finance. The second composes freely and
is unavailable to United States persons.

That is not an accident of implementation. It is a constraint, and this paper
states it as one: a token cannot simultaneously enforce holder-level restrictions
and participate in permissionless protocols. Every design decision downstream
follows from which side of that line a given instrument sits on.

We work through the consequence for a specific case --- an artificial-intelligence
equity strategy that already exists as a model portfolio --- and conclude that the
right answer is three separate instruments rather than one clever one: a fund for
investors, a network token for the protocol, and a treasury for demonstrable
assets. We give the registration path, the filing calendar and the operating cost
of each, and we identify which of the three is unnecessary for most of what people
want from the other two.
\end{abstract}


\newpage
\tableofcontents
\newpage

\section{The strategy already exists}

The subject is concrete. An artificial-intelligence sleeve exists as model output
over real market data: eight positions with model weights, carried alongside nine
other asset-class sleeves in the same book.

\begin{table}[h]
\centering
\begin{tabular}{llrl}
\toprule
Position & Name & Model weight & Tokenized equivalent \\
\midrule
MSFT & Microsoft & 19.1\% & available \\
AVGO & Broadcom & 14.3\% & available \\
AMD & AMD & 13.7\% & available \\
SMH & Semiconductor ETF & 13.7\% & ETF class available \\
AMZN & Amazon & 12.9\% & available \\
NVDA & NVIDIA & 10.5\% & available \\
GOOGL & Alphabet & 8.9\% & available \\
META & Meta & 6.7\% & available \\
\bottomrule
\end{tabular}
\caption{The modelled sleeve. Weights are model output on real market data, not
realized fund returns, holdings or net asset value.}
\end{table}

The relevant fact is the right-hand column. A tokenized-equity programme that
launched in 2025 with roughly sixty tickers has passed five hundred assets and
\$35 billion of cumulative volume, covering United States stocks and exchange-traded
funds one-to-one, settling on a public chain rather than through the conventional
clearing system, trading continuously. Essentially every position in the sleeve
has an on-chain representation today.

So the portfolio construction question is answered before we start. What remains
is the wrapper, and the wrapper is where the entire difficulty lives.

\section{The constraint}

\begin{definition}[Holder-restricted token]
A token whose transfer function consults an allowlist, a jurisdiction rule, or a
holder registry, and reverts when the recipient fails it.
\end{definition}

\begin{definition}[Permissionless composability]
The property that an arbitrary smart contract --- an automated market maker, a
lending pool, a vault --- can hold and move a token without prior authorisation.
\end{definition}

\begin{proposition}[No token has both]
Holder restriction and permissionless composability are mutually exclusive. A
protocol contract is an address that no registry has vetted; a transfer to it
either succeeds, in which case the restriction is not enforced, or reverts, in
which case the token is not composable.
\end{proposition}

This is not a limitation anyone has failed to engineer around. It is what
enforcement means. The two live precedents sit on opposite sides of it and each
is coherent on its own terms.

\begin{table}[h]
\centering
\begin{tabularx}{\textwidth}{lXX}
\toprule
 & Registered fund token & Tokenized equity \\
\midrule
Token standard & restriction-enforcing & plain transferable \\
Compliance point & every transfer & issuance and redemption only \\
Offerable to US retail & yes & no \\
Usable in a lending pool & no & yes \\
Bridgeable to another chain & only to permissioned venues & yes \\
Registration & investment company & none by the composer \\
\bottomrule
\end{tabularx}
\caption{The two coherent designs. There is no third column.}
\end{table}

The practical corollary is worth stating bluntly because it is the mistake most
plans make: \emph{a token that is a registered fund share cannot be liquidity for
anything}. Any design that requires both properties requires two tokens.

\section{The instrument menu}

Four doors, with what each costs and what each forbids. Figures are planning
ranges for a fund of the size contemplated and require confirmation by counsel and
by the service providers quoted.

\begin{table}[h]
\centering
\begin{tabularx}{\textwidth}{lrrXX}
\toprule
Vehicle & Launch & Annual & Reaches & Composable \\
\midrule
Registered interval fund & \$250--500k & \$700k--1.5M & US retail & no \\
Private fund, accredited & \$75--200k & \$150--400k & accredited US & restricted \\
Offshore fund, non-US & \$50--150k & \$100--250k & non-US & yes \\
Treasury, no fund & \$25--75k & \$50--150k & n/a & yes \\
\bottomrule
\end{tabularx}
\caption{Vehicle cost and reach. Annual figures are recurring operating cost
before any investment is made.}
\end{table}

\subsection{What the registered path actually carries}

The heaviest door is worth itemising, because its cost is not a fee schedule but
an institution. A registered closed-end fund operating as an interval fund
requires a board with a majority of independent directors, a qualified custodian,
a chief compliance officer and a written compliance programme, an annual audit by
a registered public accounting firm, portfolio reporting monthly, an annual
census filing, a separate filing for every periodic repurchase offer, a transfer
agent, and post-effective amendments to keep the offering current.

It also imports a prohibition that is decisive here: a registered fund may not
transact with its own affiliates as principal without an exemptive order. If the
fund's sponsor also issues a network token, the fund cannot buy that token and
that token cannot be a claim on the fund. Any structure of the form \emph{our
token is backed by our fund} terminates at this sentence.

The live precedent for the registered path has been operating for six years on a
portfolio of short-duration government securities and a money market fund, and is
currently offering to repurchase six percent of its shares at net asset value.
That is the structure functioning correctly. It is also a good measure of how much
machinery is required to run the simplest possible portfolio inside it.

\section{Three instruments, not one}

The resolution is to stop trying to make one token do incompatible jobs.

\paragraph{The fund.} An investable artificial-intelligence strategy, offered to
investors, holding tokenized equities and liquid digital assets. Its token is
restricted, its holders are known, and it is not composable. This is the
instrument that takes outside capital.

\paragraph{The network token.} Governance, staking and access for the protocol.
Fixed supply, plainly transferable, bridgeable, poolable, usable as liquidity in
another protocol. It is \emph{not} a claim on the fund, and the difference between
a claim and an association is the entire reason it can remain unrestricted.

\paragraph{The treasury.} Assets held by the foundation, on-chain, at published
addresses, verifiable by anyone. Not a pooled vehicle, not offered to investors,
therefore not a fund and not registrable as one.

The treasury is the instrument most plans overlook, and for the specific goal of
\emph{demonstrable backing} it dominates the other two. A verifiable treasury is
not a claim: holders can confirm the assets exist without holding a redeemable
interest in them. That distinction is what keeps the network token unrestricted,
and it is why the language used to describe it is load-bearing rather than
cosmetic. Publishing an address and an amount is a statement of fact. Promising
that a token is backed by that amount is an undertaking, and an undertaking of
returns from a common enterprise is the classic securities test reciting itself.

\begin{table}[h]
\centering
\begin{tabularx}{\textwidth}{lXXX}
\toprule
 & Fund & Network token & Treasury \\
\midrule
Purpose & investable strategy & protocol utility & demonstrable assets \\
Holders & known, restricted & anyone & n/a \\
Composable & no & yes & yes \\
Redeemable & yes, periodic & no & no \\
Registration & by vehicle chosen & exemption or safe harbour & none \\
Annual cost & \$100k--1.5M & minimal & \$50--150k \\
\bottomrule
\end{tabularx}
\caption{Division of labour. The relationships between the three are commercial,
never collateral.}
\end{table}

\section{Filing map}

What each path requires, in the order it is required.

\begin{table}[h]
\centering
\begin{tabularx}{\textwidth}{llX}
\toprule
Path & Principal filing & Ongoing \\
\midrule
Registered interval fund & registration statement for a closed-end fund
 & post-effective amendments; monthly portfolio reports; annual census; a filing
   per repurchase offer; annual audited financials \\
Private fund (accredited) & exempt-offering notice
 & state notice filings; adviser registration or exempt-reporting-adviser filing;
   annual audited financials to investors \\
Offshore fund (non-US) & none with the Commission
 & fund-domicile reporting; administrator and audit per offering documents \\
Network token & none today
 & if the proposed crypto-asset regulation is adopted: audited financials and
   periodic reports at the higher tier, and a certification to exit the
   investment-contract characterisation \\
Treasury & none & published addresses; periodic attestation \\
\bottomrule
\end{tabularx}
\caption{Filing obligations by path. Form numbers deliberately omitted; the
proposed crypto-asset regulation is a proposal and confers nothing until adopted.}
\end{table}

Two sequencing notes. A registration statement removes the offering size limit
entirely, so an exempt raise cap is irrelevant on that path and stacking the two
for the same instrument is incoherent. And the exempt path for a non-security
digital asset cannot carry an interest in a portfolio of securities, because its
own predicate is that the asset is not itself a security --- which a claim on a
securities portfolio is.

\section{On-chain architecture}

For whichever vehicle ends up composable, the engineering is the same and is
independent of the filing.

\paragraph{Vault.} A tokenized vault holding three sleeves --- tokenized equities
tracking the model, liquid digital assets, and a rebalancing buffer --- with
in-kind creation and redemption for authorised participants and buffer-denominated
subscription for everyone else.

\paragraph{Valuation across closed markets.} The hard problem, and the one most
implementations get wrong. Tokenized equities trade continuously; their underlying
shares do not. Between Friday's close and Monday's open the token price is a
forecast rather than a valuation. The vault's oracle must carry an explicit stance
for closed and halted markets rather than a stale last price, and any protocol
using vault shares as collateral must widen its margin over exactly those windows.

\paragraph{Issuer concentration.} A single issuer stands behind the entire
tokenized-equity sleeve, holding the underlying shares through one legal entity.
That is a concentration to diversify across issuers as alternatives mature, not to
assume away.

\paragraph{Redemption asymmetry.} The vault composes freely; converting its
holdings back into registered shares still passes through the issuer's compliance
gate. Anyone modelling the position as fully liquid to cash is mistaken, and the
disclosure should say so.

\section{What we do not claim}

Performance figures associated with this strategy are asset-class reference data
or model output on real market data. They are not realized fund returns, net asset
value or assets under management, and nothing here should be read as reporting
any. The vehicles described are not offered by this document; any offering is made
only by definitive documents to eligible investors. The cost ranges are planning
estimates. The proposed crypto-asset regulation referred to above is a proposal
open for comment and confers no exemption on anyone at the time of writing.

The structural conclusions --- that holder restriction and composability cannot
coexist in one token, that a registered fund cannot hold its affiliate's token,
and that a verifiable treasury is not a claim --- are the load-bearing ones, and
each should be confirmed by counsel before it is relied upon.

\end{document}
